\section{Causal Scorpion Detection}
\label{sec:scorpion}

The foundational paper's Proposition~2 \citep{lasser2026gfm} detects
scorpions through two channels: contraction attribution (correlational) and
deception detection (behavioral trust decay). This section upgrades
channel~(a) from correlational to causal, using the SCM of
Section~\ref{sec:causal}, and provides bounds on non-stationary scorpion
evasion.

\subsection{Causal Detection}

\begin{proposition}[Causal Scorpion Detection]
\label{prop:causal_scorpion}
Under the SCM $\SCM$ (Definition~\ref{def:scm}), a GFM actor can detect
    scorpion agents through two upgraded channels:
\begin{enumerate}
\item[\textbf{(a)}] \textbf{Causal contraction attribution.}
    $\CausalAttr(a_j, t)$ (Definition~\ref{def:causal_attribution}) replaces
    the correlational $\Deltahat$ signal from Proposition~2(a) of
    \citet{lasser2026gfm}. Within each stationary epoch of $a_j$'s
    strategy, the causal-attribution sequence
    $\CausalAttr(a_j, 1), \CausalAttr(a_j, 2), \ldots$ is i.i.d.\ with
    unknown mean $\mu_j$ and variance $\sigma_{\mathrm{do}}^2$. The
    actor maintains two auxiliary EWMA statistics --- a
    \emph{signed-mean} estimator and a \emph{variance} estimator ---
    for monitoring:
\begin{align}
\label{eq:causal_trust_mean}
\mu_j^{\mathrm{causal}}(t) &= \alpha \cdot \mu_j^{\mathrm{causal}}(t-1)
    + (1 - \alpha) \cdot \CausalAttr(a_j, t) \\
\label{eq:causal_trust_var}
\sigma_j^{\mathrm{causal},2}(t) &= \alpha \cdot
    \sigma_j^{\mathrm{causal},2}(t-1)
    + (1 - \alpha) \cdot
    \bigl(\CausalAttr(a_j, t) - \mu_j^{\mathrm{causal}}(t)\bigr)^2
\end{align}
    These statistics satisfy the same $L^2$ convergence guarantees as
    Proposition~\ref{prop:risk_convergence}: asymptotically unbiased
    with stationary variance bounded by
    $\tfrac{1 - \alpha}{1 + \alpha} \sigma_{\mathrm{do}}^2$ for the
    signed mean. They do not converge almost surely to a deterministic
    limit, and we deliberately do not use them as the flag condition.

    The flag condition is instead a sequential probability ratio test
    (SPRT) \citep{wald1945sequential} on the causal-attribution
    sequence, with null hypothesis
    $H_0: \mu_j \geq 0$ (non-scorpion) versus alternative
    $H_1: \mu_j \leq -\epsilon$ (scorpion with minimum detectable
    contraction magnitude $\epsilon > 0$). The actor accumulates the
    log-likelihood ratio
    $\Lambda_j(t) = \sum_{s=1}^{t} \log\bigl(p_{H_1}(\CausalAttr(a_j,
    s)) / p_{H_0}(\CausalAttr(a_j, s))\bigr)$ and flags $a_j$ as a
    scorpion when $\Lambda_j(t) > \log((1 - \beta_{\mathrm{err}}) /
    \alpha_{\mathrm{err}})$, where $\alpha_{\mathrm{err}}$ and
    $\beta_{\mathrm{err}}$ are the target Type~I and Type~II error
    rates. Under this decision rule the expected sample size is
    $\mathbb{E}[\tau_{\mathrm{detect}}] =
    O(\sigma_{\mathrm{do}}^2 / \mu_j^2)$ by Wald's identity, strictly
    fewer than the correlational detector when confounders are
    present (Proposition~\ref{prop:causal_dominates}). The signed-mean
    EWMA is retained as a human-interpretable monitor of the
    estimator state; the SPRT statistic $\Lambda_j(t)$ is the
    load-bearing decision variable.
\item[\textbf{(b)}] \textbf{Deception detection (unchanged).} The
    behavioral trust factor $T_j$ detects agents whose reported capability
    changes diverge from observations, through the same prediction-residual
    dynamics as the foundational paper (Equations~25, 33, 38 of
    \citet{lasser2026gfm}). Causal attribution does not affect this channel:
    a deceptive agent accumulates prediction residuals regardless of whether
    the contraction is causally attributed.
\end{enumerate}
\end{proposition}

\begin{proof}
Channel~(a): By the SCM, $\CausalAttr(a_j, t)$ isolates $a_j$'s causal
contribution from confounders
(Definition~\ref{def:causal_attribution}). Within a stationary epoch
of $a_j$'s strategy, $\{\CausalAttr(a_j, t)\}$ is i.i.d.\ with mean
$\mu_j$ and variance $\sigma_{\mathrm{do}}^2$; for a scorpion
$\mu_j \leq -\epsilon < 0$. The auxiliary EWMAs
$\mu_j^{\mathrm{causal}}(t)$ and $\sigma_j^{\mathrm{causal},2}(t)$
inherit the $L^2$ guarantees of
Proposition~\ref{prop:risk_convergence}: asymptotically unbiased
estimators of $\mu_j$ and $\sigma_{\mathrm{do}}^2$ with bounded
stationary variance. They are not used for the detection decision,
so we do not require almost-sure convergence.

The SPRT decision on the causal-attribution sequence is a standard
sequential test: under $H_0$ versus $H_1$ with bounded
log-likelihood ratios and finite Kullback-Leibler divergence
$D(H_1 \,\|\, H_0) > 0$, Wald's identity gives expected stopping time
$\mathbb{E}[\tau_{\mathrm{detect}}] \approx
\log((1 - \beta_{\mathrm{err}}) / \alpha_{\mathrm{err}}) /
D(H_1 \,\|\, H_0)$. For Gaussian causal attributions with
$D(H_1 \,\|\, H_0) = \epsilon^2 / (2 \sigma_{\mathrm{do}}^2)$, this
reduces to $\mathbb{E}[\tau_{\mathrm{detect}}] =
O(\sigma_{\mathrm{do}}^2 / \epsilon^2)$; substituting the actual mean
$\mu_j$ (which lower-bounds $\epsilon$ for a detectable scorpion)
gives $O(\sigma_{\mathrm{do}}^2 / \mu_j^2)$, matching the rate claimed
in Proposition~\ref{prop:causal_dominates}(b). The SPRT Type~I error
is bounded by $\alpha_{\mathrm{err}}$ and Type~II by
$\beta_{\mathrm{err}}$ by construction.

Channel~(b): The behavioral trust dynamics (Appendix~B of
\citet{lasser2026gfm}) operate on behavioral prediction residuals
$r_j(t) = R_j(t) - \hat{R}_j(t)$, which are independent of causal
attribution. A deceptive scorpion that misreports capability changes
accumulates $r_j(t) \neq 0$ regardless of confounders, driving
$T_j \to T_j^* < 1$.
\end{proof}

\subsection{Non-Stationary Scorpion Bounds}

The foundational paper notes that ``a non-stationary scorpion that adapts its strategy in
response to detection pressure can evade both channels indefinitely''
(Proposition~2 of \citet{lasser2026gfm}). Under the SCM, we can bound the
adaptation rate required for evasion.

\begin{proposition}[Evasion Bandwidth]
\label{prop:evasion_bandwidth}
A scorpion $a_j$ that changes strategy every $\tau_{\mathrm{adapt}}$ steps
    produces at most $\tau_{\mathrm{adapt}}$ observations per strategy
    epoch. The causal detector requires
    $\tau_{\mathrm{detect}} = O(\sigma_{\mathrm{do}}^2 / \mu_j^2)$
    observations per epoch to identify the scorpion. If
    $\tau_{\mathrm{adapt}} > \tau_{\mathrm{detect}}$, the scorpion is
    detected within each epoch. If $\tau_{\mathrm{adapt}} <
    \tau_{\mathrm{detect}}$, the scorpion evades detection by changing
    strategy before enough evidence accumulates.
\end{proposition}

\begin{proof}
Within each epoch of length $\tau_{\mathrm{adapt}}$, the scorpion's
strategy is stationary and the causal-attribution sequence is i.i.d.\
with mean $\mu_j$ and variance $\sigma_{\mathrm{do}}^2$. The SPRT
statistic accumulates log-likelihood-ratio evidence at expected rate
$D(H_1 \,\|\, H_0) = \mu_j^2 / (2 \sigma_{\mathrm{do}}^2)$ per
observation under the Gaussian surrogate used in
Proposition~\ref{prop:causal_scorpion}. After $\tau_{\mathrm{adapt}}$
observations, the expected accumulated evidence is
$\tau_{\mathrm{adapt}} \cdot \mu_j^2 / (2 \sigma_{\mathrm{do}}^2)$. An
epoch boundary resets the SPRT statistic (the null hypothesis of the
test is epoch-local). If $\tau_{\mathrm{adapt}} <
\tau_{\mathrm{detect}}$, the expected evidence at the epoch boundary
is below the SPRT acceptance boundary
$\log((1 - \beta_{\mathrm{err}}) / \alpha_{\mathrm{err}})$ and the
scorpion evades detection within that epoch.
\end{proof}

\paragraph{Interpretation.} The evasion bandwidth
$\tau_{\mathrm{adapt}}^* = O(\sigma_{\mathrm{do}}^2 / \mu_j^2)$ is the
minimum adaptation rate a scorpion needs to evade the causal detector.
A fast-adapting scorpion ($\tau_{\mathrm{adapt}} \ll
\tau_{\mathrm{adapt}}^*$) evades but is constrained to small contractions
per epoch (otherwise $\mu_j$ is large and $\tau_{\mathrm{adapt}}^*$ shrinks).
A slow-adapting scorpion ($\tau_{\mathrm{adapt}} \gg
\tau_{\mathrm{adapt}}^*$) is detected within each epoch. The product
$\mu_j \cdot \tau_{\mathrm{adapt}}$ (total contraction per epoch) is
bounded: a scorpion can either cause large damage slowly (and be detected)
or small damage quickly (and evade), but not large damage quickly.

The time-averaged undetected contraction rate for an evading scorpion
is bounded by the signed mean $|\mu_j|$ itself, not by an independent
$\sigma_{\mathrm{do}}^2 / |\mu_j|$ factor: the evasion constraint
$\tau_{\mathrm{adapt}} < \tau_{\mathrm{detect}} =
\sigma_{\mathrm{do}}^2 / \mu_j^2$ (ignoring the constant
$\log(1/\alpha_{\mathrm{err}})$) bounds the epoch length, and within
each epoch the per-step contraction is $|\mu_j|$. The appendix proof
(§\ref{app:proof_evasion}) establishes this bound via the sequential
probability ratio test accumulation rate. Substituting the evasion
constraint $|\mu_j| < \sigma_{\mathrm{do}} / \sqrt{\tau_{\mathrm{adapt}}}$
gives an equivalent expression in terms of the adaptation budget:
\begin{equation}
\label{eq:evasion_bound}
|\text{undetected contraction per step}| \;\leq\; |\mu_j|
    \;\leq\; \frac{\sigma_{\mathrm{do}}}{\sqrt{\tau_{\mathrm{adapt}}}}
\end{equation}
The bound tightens as the causal model improves ($\sigma_{\mathrm{do}}$
decreases) and as the scorpion's adaptation budget
$\tau_{\mathrm{adapt}}$ grows: a scorpion that adapts frequently (small
$\tau_{\mathrm{adapt}}$) can only afford very small per-step
contractions, while a scorpion that adapts slowly can sustain larger
contractions per step but for correspondingly fewer steps before
detection.
